%
% ei.tex -- draw Ei graph
%
% (c) 2021 Prof Dr Andreas Müller, OST Ostschweizer Fachhochschule
%
\documentclass[tikz]{standalone}
\usepackage{amsmath}
\usepackage{times}
\usepackage{txfonts}
\usepackage{pgfplots}
\usepackage{csvsimple}
\usetikzlibrary{arrows,intersections,math,calc}
\definecolor{darkred}{rgb}{0.8,0,0}
\definecolor{darkgreen}{rgb}{0,0.6,0}
\begin{document}
\def\skala{2}
\def\dx{1}
\def\dy{0.25}
\def\kurvelinks{
	({0.0010*\dx},{-6.3295*\dy})
	--({0.0166*\dx},{-3.5061*\dy})
	--({0.0321*\dx},{-2.8279*\dy})
	--({0.0477*\dx},{-2.4169*\dy})
	--({0.0633*\dx},{-2.1185*\dy})
	--({0.0789*\dx},{-1.8824*\dy})
	--({0.0944*\dx},{-1.6859*\dy})
	--({0.1100*\dx},{-1.5169*\dy})
	--({0.1256*\dx},{-1.3679*\dy})
	--({0.1412*\dx},{-1.2344*\dy})
	--({0.1567*\dx},{-1.1129*\dy})
	--({0.1723*\dx},{-1.0013*\dy})
	--({0.1879*\dx},{-0.8977*\dy})
	--({0.2034*\dx},{-0.8009*\dy})
	--({0.2190*\dx},{-0.7098*\dy})
	--({0.2346*\dx},{-0.6236*\dy})
	--({0.2502*\dx},{-0.5417*\dy})
	--({0.2657*\dx},{-0.4635*\dy})
	--({0.2813*\dx},{-0.3887*\dy})
	--({0.2969*\dx},{-0.3167*\dy})
	--({0.3125*\dx},{-0.2474*\dy})
	--({0.3280*\dx},{-0.1804*\dy})
	--({0.3436*\dx},{-0.1155*\dy})
	--({0.3592*\dx},{-0.0525*\dy})
	--({0.3747*\dx},{0.0087*\dy})
	--({0.3903*\dx},{0.0684*\dy})
	--({0.4059*\dx},{0.1267*\dy})
	--({0.4215*\dx},{0.1836*\dy})
	--({0.4370*\dx},{0.2393*\dy})
	--({0.4526*\dx},{0.2939*\dy})
	--({0.4682*\dx},{0.3476*\dy})
	--({0.4838*\dx},{0.4002*\dy})
	--({0.4993*\dx},{0.4520*\dy})
	--({0.5149*\dx},{0.5030*\dy})
	--({0.5305*\dx},{0.5533*\dy})
	--({0.5461*\dx},{0.6028*\dy})
	--({0.5616*\dx},{0.6517*\dy})
	--({0.5772*\dx},{0.7001*\dy})
	--({0.5928*\dx},{0.7479*\dy})
	--({0.6083*\dx},{0.7951*\dy})
	--({0.6239*\dx},{0.8419*\dy})
	--({0.6395*\dx},{0.8883*\dy})
	--({0.6551*\dx},{0.9343*\dy})
	--({0.6706*\dx},{0.9799*\dy})
	--({0.6862*\dx},{1.0251*\dy})
	--({0.7018*\dx},{1.0700*\dy})
	--({0.7174*\dx},{1.1146*\dy})
	--({0.7329*\dx},{1.1590*\dy})
	--({0.7485*\dx},{1.2031*\dy})
	--({0.7641*\dx},{1.2470*\dy})
	--({0.7796*\dx},{1.2906*\dy})
	--({0.7952*\dx},{1.3341*\dy})
	--({0.8108*\dx},{1.3774*\dy})
	--({0.8264*\dx},{1.4205*\dy})
	--({0.8419*\dx},{1.4635*\dy})
	--({0.8575*\dx},{1.5064*\dy})
	--({0.8731*\dx},{1.5491*\dy})
	--({0.8887*\dx},{1.5918*\dy})
	--({0.9042*\dx},{1.6344*\dy})
	--({0.9198*\dx},{1.6769*\dy})
	--({0.9354*\dx},{1.7193*\dy})
	--({0.9509*\dx},{1.7617*\dy})
	--({0.9665*\dx},{1.8041*\dy})
	--({0.9821*\dx},{1.8464*\dy})
	--({0.9977*\dx},{1.8888*\dy})
	--({1.0132*\dx},{1.9311*\dy})
	--({1.0288*\dx},{1.9734*\dy})
	--({1.0444*\dx},{2.0158*\dy})
	--({1.0600*\dx},{2.0582*\dy})
	--({1.0755*\dx},{2.1006*\dy})
	--({1.0911*\dx},{2.1431*\dy})
	--({1.1067*\dx},{2.1856*\dy})
	--({1.1222*\dx},{2.2282*\dy})
	--({1.1378*\dx},{2.2709*\dy})
	--({1.1534*\dx},{2.3136*\dy})
	--({1.1690*\dx},{2.3564*\dy})
	--({1.1845*\dx},{2.3994*\dy})
	--({1.2001*\dx},{2.4424*\dy})
	--({1.2157*\dx},{2.4855*\dy})
	--({1.2313*\dx},{2.5288*\dy})
	--({1.2468*\dx},{2.5722*\dy})
	--({1.2624*\dx},{2.6157*\dy})
	--({1.2780*\dx},{2.6594*\dy})
	--({1.2935*\dx},{2.7032*\dy})
	--({1.3091*\dx},{2.7472*\dy})
	--({1.3247*\dx},{2.7913*\dy})
	--({1.3403*\dx},{2.8356*\dy})
	--({1.3558*\dx},{2.8801*\dy})
	--({1.3714*\dx},{2.9247*\dy})
	--({1.3870*\dx},{2.9696*\dy})
	--({1.4026*\dx},{3.0146*\dy})
	--({1.4181*\dx},{3.0599*\dy})
	--({1.4337*\dx},{3.1053*\dy})
	--({1.4493*\dx},{3.1510*\dy})
	--({1.4648*\dx},{3.1969*\dy})
	--({1.4804*\dx},{3.2430*\dy})
	--({1.4960*\dx},{3.2893*\dy})
	--({1.5116*\dx},{3.3359*\dy})
	--({1.5271*\dx},{3.3827*\dy})
	--({1.5427*\dx},{3.4298*\dy})
	--({1.5583*\dx},{3.4772*\dy})
	--({1.5739*\dx},{3.5248*\dy})
	--({1.5894*\dx},{3.5727*\dy})
	--({1.6050*\dx},{3.6208*\dy})
	--({1.6206*\dx},{3.6693*\dy})
	--({1.6362*\dx},{3.7180*\dy})
	--({1.6517*\dx},{3.7670*\dy})
	--({1.6673*\dx},{3.8164*\dy})
	--({1.6829*\dx},{3.8660*\dy})
	--({1.6984*\dx},{3.9159*\dy})
	--({1.7140*\dx},{3.9662*\dy})
	--({1.7296*\dx},{4.0168*\dy})
	--({1.7452*\dx},{4.0678*\dy})
	--({1.7607*\dx},{4.1190*\dy})
	--({1.7763*\dx},{4.1707*\dy})
	--({1.7919*\dx},{4.2226*\dy})
	--({1.8075*\dx},{4.2750*\dy})
	--({1.8230*\dx},{4.3277*\dy})
	--({1.8386*\dx},{4.3807*\dy})
	--({1.8542*\dx},{4.4342*\dy})
	--({1.8697*\dx},{4.4880*\dy})
	--({1.8853*\dx},{4.5422*\dy})
	--({1.9009*\dx},{4.5968*\dy})
	--({1.9165*\dx},{4.6519*\dy})
	--({1.9320*\dx},{4.7073*\dy})
	--({1.9476*\dx},{4.7632*\dy})
	--({1.9632*\dx},{4.8194*\dy})
	--({1.9788*\dx},{4.8762*\dy})
	--({1.9943*\dx},{4.9333*\dy})
	--({2.0099*\dx},{4.9909*\dy})
	--({2.0255*\dx},{5.0489*\dy})
	--({2.0410*\dx},{5.1075*\dy})
	--({2.0566*\dx},{5.1664*\dy})
	--({2.0722*\dx},{5.2259*\dy})
	--({2.0878*\dx},{5.2858*\dy})
	--({2.1033*\dx},{5.3462*\dy})
	--({2.1189*\dx},{5.4071*\dy})
	--({2.1345*\dx},{5.4685*\dy})
	--({2.1501*\dx},{5.5305*\dy})
	--({2.1656*\dx},{5.5929*\dy})
	--({2.1812*\dx},{5.6559*\dy})
	--({2.1968*\dx},{5.7194*\dy})
	--({2.2123*\dx},{5.7834*\dy})
	--({2.2279*\dx},{5.8480*\dy})
	--({2.2435*\dx},{5.9132*\dy})
	--({2.2591*\dx},{5.9789*\dy})
	--({2.2746*\dx},{6.0452*\dy})
	--({2.2902*\dx},{6.1120*\dy})
	--({2.3058*\dx},{6.1795*\dy})
	--({2.3214*\dx},{6.2476*\dy})
	--({2.3369*\dx},{6.3162*\dy})
	--({2.3525*\dx},{6.3855*\dy})
	--({2.3681*\dx},{6.4554*\dy})
	--({2.3836*\dx},{6.5259*\dy})
	--({2.3992*\dx},{6.5971*\dy})
	--({2.4148*\dx},{6.6689*\dy})
	--({2.4304*\dx},{6.7414*\dy})
	--({2.4459*\dx},{6.8145*\dy})
	--({2.4615*\dx},{6.8884*\dy})
	--({2.4771*\dx},{6.9629*\dy})
	--({2.4927*\dx},{7.0381*\dy})
	--({2.5082*\dx},{7.1140*\dy})
	--({2.5238*\dx},{7.1906*\dy})
	--({2.5394*\dx},{7.2679*\dy})
	--({2.5549*\dx},{7.3460*\dy})
	--({2.5705*\dx},{7.4248*\dy})
	--({2.5861*\dx},{7.5044*\dy})
	--({2.6017*\dx},{7.5848*\dy})
	--({2.6172*\dx},{7.6659*\dy})
	--({2.6328*\dx},{7.7478*\dy})
	--({2.6484*\dx},{7.8305*\dy})
	--({2.6640*\dx},{7.9140*\dy})
	--({2.6795*\dx},{7.9983*\dy})
	--({2.6951*\dx},{8.0834*\dy})
	--({2.7107*\dx},{8.1694*\dy})
	--({2.7263*\dx},{8.2562*\dy})
	--({2.7418*\dx},{8.3439*\dy})
	--({2.7574*\dx},{8.4325*\dy})
	--({2.7730*\dx},{8.5219*\dy})
	--({2.7885*\dx},{8.6123*\dy})
	--({2.8041*\dx},{8.7035*\dy})
	--({2.8197*\dx},{8.7957*\dy})
	--({2.8353*\dx},{8.8888*\dy})
	--({2.8508*\dx},{8.9828*\dy})
	--({2.8664*\dx},{9.0778*\dy})
	--({2.8820*\dx},{9.1738*\dy})
	--({2.8976*\dx},{9.2707*\dy})
	--({2.9131*\dx},{9.3686*\dy})
	--({2.9287*\dx},{9.4676*\dy})
	--({2.9443*\dx},{9.5675*\dy})
	--({2.9598*\dx},{9.6685*\dy})
	--({2.9754*\dx},{9.7706*\dy})
	--({2.9910*\dx},{9.8737*\dy})
	--({3.0066*\dx},{9.9779*\dy})
	--({3.0221*\dx},{10.0831*\dy})
	--({3.0377*\dx},{10.1895*\dy})
	--({3.0533*\dx},{10.2970*\dy})
	--({3.0689*\dx},{10.4056*\dy})
	--({3.0844*\dx},{10.5154*\dy})
	--({3.1000*\dx},{10.6263*\dy})
}
\def\kurverechts{
	({-0.0010*\dx},{-6.3315*\dy})
	--({-0.0166*\dx},{-3.5393*\dy})
	--({-0.0321*\dx},{-2.8921*\dy})
	--({-0.0477*\dx},{-2.5124*\dy})
	--({-0.0633*\dx},{-2.2451*\dy})
	--({-0.0789*\dx},{-2.0401*\dy})
	--({-0.0944*\dx},{-1.8749*\dy})
	--({-0.1100*\dx},{-1.7370*\dy})
	--({-0.1256*\dx},{-1.6193*\dy})
	--({-0.1412*\dx},{-1.5170*\dy})
	--({-0.1567*\dx},{-1.4268*\dy})
	--({-0.1723*\dx},{-1.3464*\dy})
	--({-0.1879*\dx},{-1.2742*\dy})
	--({-0.2034*\dx},{-1.2087*\dy})
	--({-0.2190*\dx},{-1.1490*\dy})
	--({-0.2346*\dx},{-1.0942*\dy})
	--({-0.2502*\dx},{-1.0438*\dy})
	--({-0.2657*\dx},{-0.9971*\dy})
	--({-0.2813*\dx},{-0.9538*\dy})
	--({-0.2969*\dx},{-0.9134*\dy})
	--({-0.3125*\dx},{-0.8757*\dy})
	--({-0.3280*\dx},{-0.8404*\dy})
	--({-0.3436*\dx},{-0.8073*\dy})
	--({-0.3592*\dx},{-0.7761*\dy})
	--({-0.3747*\dx},{-0.7467*\dy})
	--({-0.3903*\dx},{-0.7189*\dy})
	--({-0.4059*\dx},{-0.6926*\dy})
	--({-0.4215*\dx},{-0.6677*\dy})
	--({-0.4370*\dx},{-0.6441*\dy})
	--({-0.4526*\dx},{-0.6216*\dy})
	--({-0.4682*\dx},{-0.6003*\dy})
	--({-0.4838*\dx},{-0.5800*\dy})
	--({-0.4993*\dx},{-0.5606*\dy})
	--({-0.5149*\dx},{-0.5421*\dy})
	--({-0.5305*\dx},{-0.5244*\dy})
	--({-0.5461*\dx},{-0.5075*\dy})
	--({-0.5616*\dx},{-0.4914*\dy})
	--({-0.5772*\dx},{-0.4759*\dy})
	--({-0.5928*\dx},{-0.4611*\dy})
	--({-0.6083*\dx},{-0.4468*\dy})
	--({-0.6239*\dx},{-0.4332*\dy})
	--({-0.6395*\dx},{-0.4201*\dy})
	--({-0.6551*\dx},{-0.4075*\dy})
	--({-0.6706*\dx},{-0.3954*\dy})
	--({-0.6862*\dx},{-0.3837*\dy})
	--({-0.7018*\dx},{-0.3725*\dy})
	--({-0.7174*\dx},{-0.3617*\dy})
	--({-0.7329*\dx},{-0.3513*\dy})
	--({-0.7485*\dx},{-0.3413*\dy})
	--({-0.7641*\dx},{-0.3316*\dy})
	--({-0.7796*\dx},{-0.3223*\dy})
	--({-0.7952*\dx},{-0.3133*\dy})
	--({-0.8108*\dx},{-0.3046*\dy})
	--({-0.8264*\dx},{-0.2962*\dy})
	--({-0.8419*\dx},{-0.2881*\dy})
	--({-0.8575*\dx},{-0.2803*\dy})
	--({-0.8731*\dx},{-0.2727*\dy})
	--({-0.8887*\dx},{-0.2654*\dy})
	--({-0.9042*\dx},{-0.2583*\dy})
	--({-0.9198*\dx},{-0.2514*\dy})
	--({-0.9354*\dx},{-0.2448*\dy})
	--({-0.9509*\dx},{-0.2384*\dy})
	--({-0.9665*\dx},{-0.2321*\dy})
	--({-0.9821*\dx},{-0.2261*\dy})
	--({-0.9977*\dx},{-0.2202*\dy})
	--({-1.0132*\dx},{-0.2146*\dy})
	--({-1.0288*\dx},{-0.2091*\dy})
	--({-1.0444*\dx},{-0.2038*\dy})
	--({-1.0600*\dx},{-0.1986*\dy})
	--({-1.0755*\dx},{-0.1936*\dy})
	--({-1.0911*\dx},{-0.1887*\dy})
	--({-1.1067*\dx},{-0.1840*\dy})
	--({-1.1222*\dx},{-0.1794*\dy})
	--({-1.1378*\dx},{-0.1749*\dy})
	--({-1.1534*\dx},{-0.1706*\dy})
	--({-1.1690*\dx},{-0.1664*\dy})
	--({-1.1845*\dx},{-0.1623*\dy})
	--({-1.2001*\dx},{-0.1584*\dy})
	--({-1.2157*\dx},{-0.1545*\dy})
	--({-1.2313*\dx},{-0.1508*\dy})
	--({-1.2468*\dx},{-0.1471*\dy})
	--({-1.2624*\dx},{-0.1436*\dy})
	--({-1.2780*\dx},{-0.1402*\dy})
	--({-1.2935*\dx},{-0.1368*\dy})
	--({-1.3091*\dx},{-0.1336*\dy})
	--({-1.3247*\dx},{-0.1304*\dy})
	--({-1.3403*\dx},{-0.1273*\dy})
	--({-1.3558*\dx},{-0.1243*\dy})
	--({-1.3714*\dx},{-0.1214*\dy})
	--({-1.3870*\dx},{-0.1185*\dy})
	--({-1.4026*\dx},{-0.1158*\dy})
	--({-1.4181*\dx},{-0.1131*\dy})
	--({-1.4337*\dx},{-0.1105*\dy})
	--({-1.4493*\dx},{-0.1079*\dy})
	--({-1.4648*\dx},{-0.1054*\dy})
	--({-1.4804*\dx},{-0.1030*\dy})
	--({-1.4960*\dx},{-0.1006*\dy})
	--({-1.5116*\dx},{-0.0983*\dy})
	--({-1.5271*\dx},{-0.0961*\dy})
	--({-1.5427*\dx},{-0.0939*\dy})
	--({-1.5583*\dx},{-0.0918*\dy})
	--({-1.5739*\dx},{-0.0897*\dy})
	--({-1.5894*\dx},{-0.0877*\dy})
	--({-1.6050*\dx},{-0.0857*\dy})
	--({-1.6206*\dx},{-0.0838*\dy})
	--({-1.6362*\dx},{-0.0819*\dy})
	--({-1.6517*\dx},{-0.0800*\dy})
	--({-1.6673*\dx},{-0.0783*\dy})
	--({-1.6829*\dx},{-0.0765*\dy})
	--({-1.6984*\dx},{-0.0748*\dy})
	--({-1.7140*\dx},{-0.0732*\dy})
	--({-1.7296*\dx},{-0.0715*\dy})
	--({-1.7452*\dx},{-0.0700*\dy})
	--({-1.7607*\dx},{-0.0684*\dy})
	--({-1.7763*\dx},{-0.0669*\dy})
	--({-1.7919*\dx},{-0.0655*\dy})
	--({-1.8075*\dx},{-0.0640*\dy})
	--({-1.8230*\dx},{-0.0626*\dy})
	--({-1.8386*\dx},{-0.0613*\dy})
	--({-1.8542*\dx},{-0.0599*\dy})
	--({-1.8697*\dx},{-0.0586*\dy})
	--({-1.8853*\dx},{-0.0574*\dy})
	--({-1.9009*\dx},{-0.0561*\dy})
	--({-1.9165*\dx},{-0.0549*\dy})
	--({-1.9320*\dx},{-0.0537*\dy})
	--({-1.9476*\dx},{-0.0526*\dy})
	--({-1.9632*\dx},{-0.0515*\dy})
	--({-1.9788*\dx},{-0.0504*\dy})
	--({-1.9943*\dx},{-0.0493*\dy})
	--({-2.0099*\dx},{-0.0482*\dy})
	--({-2.0255*\dx},{-0.0472*\dy})
	--({-2.0410*\dx},{-0.0462*\dy})
	--({-2.0566*\dx},{-0.0452*\dy})
	--({-2.0722*\dx},{-0.0443*\dy})
	--({-2.0878*\dx},{-0.0433*\dy})
	--({-2.1033*\dx},{-0.0424*\dy})
	--({-2.1189*\dx},{-0.0415*\dy})
	--({-2.1345*\dx},{-0.0407*\dy})
	--({-2.1501*\dx},{-0.0398*\dy})
	--({-2.1656*\dx},{-0.0390*\dy})
	--({-2.1812*\dx},{-0.0382*\dy})
	--({-2.1968*\dx},{-0.0374*\dy})
	--({-2.2123*\dx},{-0.0366*\dy})
	--({-2.2279*\dx},{-0.0358*\dy})
	--({-2.2435*\dx},{-0.0351*\dy})
	--({-2.2591*\dx},{-0.0343*\dy})
	--({-2.2746*\dx},{-0.0336*\dy})
	--({-2.2902*\dx},{-0.0329*\dy})
	--({-2.3058*\dx},{-0.0323*\dy})
	--({-2.3214*\dx},{-0.0316*\dy})
	--({-2.3369*\dx},{-0.0309*\dy})
	--({-2.3525*\dx},{-0.0303*\dy})
	--({-2.3681*\dx},{-0.0297*\dy})
	--({-2.3836*\dx},{-0.0291*\dy})
	--({-2.3992*\dx},{-0.0285*\dy})
	--({-2.4148*\dx},{-0.0279*\dy})
	--({-2.4304*\dx},{-0.0273*\dy})
	--({-2.4459*\dx},{-0.0268*\dy})
	--({-2.4615*\dx},{-0.0262*\dy})
	--({-2.4771*\dx},{-0.0257*\dy})
	--({-2.4927*\dx},{-0.0252*\dy})
	--({-2.5082*\dx},{-0.0246*\dy})
	--({-2.5238*\dx},{-0.0241*\dy})
	--({-2.5394*\dx},{-0.0237*\dy})
	--({-2.5549*\dx},{-0.0232*\dy})
	--({-2.5705*\dx},{-0.0227*\dy})
	--({-2.5861*\dx},{-0.0223*\dy})
	--({-2.6017*\dx},{-0.0218*\dy})
	--({-2.6172*\dx},{-0.0214*\dy})
	--({-2.6328*\dx},{-0.0209*\dy})
	--({-2.6484*\dx},{-0.0205*\dy})
	--({-2.6640*\dx},{-0.0201*\dy})
	--({-2.6795*\dx},{-0.0197*\dy})
	--({-2.6951*\dx},{-0.0193*\dy})
	--({-2.7107*\dx},{-0.0189*\dy})
	--({-2.7263*\dx},{-0.0185*\dy})
	--({-2.7418*\dx},{-0.0182*\dy})
	--({-2.7574*\dx},{-0.0178*\dy})
	--({-2.7730*\dx},{-0.0175*\dy})
	--({-2.7885*\dx},{-0.0171*\dy})
	--({-2.8041*\dx},{-0.0168*\dy})
	--({-2.8197*\dx},{-0.0164*\dy})
	--({-2.8353*\dx},{-0.0161*\dy})
	--({-2.8508*\dx},{-0.0158*\dy})
	--({-2.8664*\dx},{-0.0155*\dy})
	--({-2.8820*\dx},{-0.0152*\dy})
	--({-2.8976*\dx},{-0.0149*\dy})
	--({-2.9131*\dx},{-0.0146*\dy})
	--({-2.9287*\dx},{-0.0143*\dy})
	--({-2.9443*\dx},{-0.0140*\dy})
	--({-2.9598*\dx},{-0.0137*\dy})
	--({-2.9754*\dx},{-0.0135*\dy})
	--({-2.9910*\dx},{-0.0132*\dy})
	--({-3.0066*\dx},{-0.0129*\dy})
	--({-3.0221*\dx},{-0.0127*\dy})
	--({-3.0377*\dx},{-0.0124*\dy})
	--({-3.0533*\dx},{-0.0122*\dy})
	--({-3.0689*\dx},{-0.0120*\dy})
	--({-3.0844*\dx},{-0.0117*\dy})
	--({-3.1000*\dx},{-0.0115*\dy})
}

\begin{tikzpicture}[>=latex,thick,scale=\skala]

\clip ({-3.1*\dx},{-4.4*\dy}) rectangle ({3.2*\dx},{10.85*\dy});

\draw[->] (-3.05,0) -- (3.15,0) coordinate[label={$x$}];
\draw[->] (0,{-4.4*\dy}) -- (0,{10.5*\dy})
	coordinate[label={right:$\operatorname{Ei}(x)$}];

\foreach \x in {-3,...,-1,1,2,3}{
	\draw (\x,{-0.05/\skala}) -- ++(0,{0.1/\skala});
}
\foreach \x in {-3,...,-1}{
	\node at (\x,0) [above] {$\x\mathstrut$};
}
\foreach \x in {1,...,3}{
	\node at (\x,0) [below] {$\x\mathstrut$};
}
\foreach \y in {-4,...,-1,1,2,...,10}{
	\draw ({-0.05/\skala},{\y*\dy}) -- ++({0.1/\skala},0);
}
\foreach \y in {-4,-2,2,4,...,10}{
	\node at (0,{\y*\dy}) [left] {$\y\mathstrut$};
}

\begin{scope}
	\clip ({-3.1*\dx},{-4.4*\dy}) rectangle ({3.1*\dx},{10*\dy});
	\draw[color=darkred,line width=1.2pt] \kurvelinks;
	\draw[color=darkred,line width=1.2pt] \kurverechts;
\end{scope}

\end{tikzpicture}
\end{document}

